% Part:sets-functions-relations
% Chapter: size-of-sets
% Section: enumerations

\documentclass[../../../include/open-logic-section]{subfiles}

\begin{document}

\olfileid{sfr}{siz}{enm}

\olsection{Opsommingen en \usetoken{S}{enumerable} verzamelingen}

\begin{editorial}
  Deze paragraaf gaat over het opsommen van verzamelingen. Hier is een
  opsomming een surjectieve functie met $\PosInt$ als domein. Alles wordt
  stap voor stap uitgelegd voor lezers die nog
  niet veel wiskunde kennen. In \olref[enm-alt]{sec} staat een kortere
  versie. Daar heet een bijectie tussen $\Nat$ (of een beginstuk
  daarvan) en de verzameling een opsomming.
\end{editorial}

\begin{explain}
We hebben al voorbeelden van verzamelingen gegeven door hun
!!{element}s op een lijst te zetten. Nu bekijken we hoe je de
!!{element}s van een verzameling kunt opsommen en wanneer dat kan,
ook als de verzameling oneindig is.
\end{explain}

\begin{defn}[Opsomming, in gewone woorden]
In gewone woorden is een \emph{opsomming} van een verzameling~$A$ een
lijst van !!{element}s van~$A$. Die lijst mag oneindig zijn, maar ieder
!!{element} van $A$ moet ergens op een eindige plaats staan. Als $A$
zo'n opsomming heeft, noemen we $A$ \emph{!!{enumerable}}.
\end{defn}

\begin{explain}
Bij zo'n opsomming moet je op een paar dingen letten:
\begin{enumerate}
\item Voor ons telt een lijst alleen als opsomming als zij een begin
  heeft. Ieder !!{element}, behalve het eerste, moet precies één
  !!{element} vlak vóór zich hebben. Er staan dus maar eindig veel
  elementen tussen het eerste !!{element} van de lijst en elk ander
  !!{element}. Elk !!{element} van een opsomming heeft daardoor een
  eindige plaats: het eerste !!{element} staat op plaats~$1$, het
  tweede op plaats~$2$, enzovoort.
\item Je kunt dezelfde verzameling~$A$ op verschillende manieren
  opsommen; alleen de volgorde van de !!{element}s is anders. De lijst
  $4$, $1$, $25$, $16$,~$9$ somt de verzameling van
  de eerste vijf kwadraten net zo goed op als $1$, $4$, $9$, $16$,~$25$.
\item Herhalingen zijn toegestaan: $1$, $1$, $2$, $2$, $3$, $3$,
  ~\dots{} somt dezelfde verzameling op als $1$, $2$, $3$,~\dots{}.
\item Maar als we één bepaalde opsomming geven, doen de volgorde en de
  herhalingen er \emph{wel} toe. We kunnen een opsomming van de positieve
  gehele getallen laten beginnen met $1$, $2$, $3$, $1$, \dots{}, maar in de gewone
  volgorde $1$, $2$, $3$, $4$,~\dots is het patroon makkelijker te zien.
\item Een opsomming moet een begin hebben. De lijst \dots, $3$, $2$,
  $1$ is dus geen opsomming van de positieve gehele getallen: zij heeft
  geen eerste !!{element}. Vraag maar op welke plaats in deze lijst het
  getal 76 staat. Dan zie je waarom dit uit de eerdere definitie volgt.
\item Deze lijst somt de positieve gehele getallen niet op: $1$, $3$,
  $5$, \dots, $2$, $4$, $6$, \dots\@ Dat werkt niet, want de even
  getallen zouden op plaats $\infty + 1$, $\infty + 2$, $\infty + 3$
  staan en niet op een eindige plaats.
\item Ook de lege verzameling is aftelbaar: de lege lijst somt haar op!{}
\end{enumerate}
\end{explain}

\begin{prop}
  Als $A$ een opsomming heeft, kan zij ook zonder herhalingen worden
  opgesomd.
\end{prop}

\begin{proof}
  Stel dat $A$ wordt opgesomd door $x_1$, $x_2$, \dots{}, waarbij elke
  $x_i$ een !!{element} van~$A$ is. We laten de latere, herhaalde !!{element}s weg.
  In de nieuwe lijst nemen we $x_i$ op als dit !!a{element} van~$A$
  niet eerder onder $x_1$, \dots,
  $x_{i-1}$ stond. We laten $x_i$ weg als het al onder $x_1$, \dots,
  ~$x_{i-1}$ stond.
\end{proof}

Om precies te kunnen bepalen hoe opsommingen en !!{enumerable}
verzamelingen werken en er iets over te kunnen bewijzen, moeten we
nauwkeuriger zeggen wat een opsomming is. Dat doen we nu.

\begin{defn}[Opsomming, precies gedefinieerd] 
Een \emph{opsomming} van een verzameling $A \neq \emptyset$ is elke
!!{surjective} functie $f \colon \PosInt \to A$. Dat betekent dat elk
element van de verzameling minstens één keer als functiewaarde voorkomt.
\end{defn}

\begin{explain}
We laten zien dat de precieze definitie en de eerdere uitleg met een
mogelijk oneindige lijst hetzelfde zeggen. Elke !!{surjective} functie
van $\PosInt$ naar een verzameling~$A$ somt~$A$ op. De functie geeft
namelijk deze lijst: $f(1)$, $f(2)$, $f(3)$, \dots. Omdat $f$
!!{surjective} is, is ieder !!{element} van~$A$ gelijk aan~$f(n)$ voor
minstens één~$n \in \PosInt$. Elk !!{element} van $A$ staat dus op een
eindige plaats in de lijst. De functie hoeft niet !!{injective} te
zijn en mag dezelfde waarde daarom vaker geven. Zoals we al zagen, zijn
zulke herhalingen toegestaan.

Omgekeerd: als een lijst alle !!{element}s van~$A$ opsomt, kunnen we een
!!a{surjective} functie $f\colon \PosInt \to A$ maken. We nemen voor
$f(n)$ het $n$-de !!{element} van de lijst. Bestaat er geen $n$-de
!!{element}, dan nemen we het laatste !!{element} van de lijst. Alleen
als $A$ leeg is, en de lijst dus ook leeg is, krijgen we zo geen
!!a{surjective} functie. Elke niet-lege lijst geeft dus een
!!a{surjective} functie $f\colon \PosInt \to A$.
\end{explain}

\begin{defn}
  \ollabel{defn:enumerable}
  Een verzameling~$A$ is !!{enumerable} als zij leeg is of een opsomming
  heeft, en alleen dan.
\end{defn}

\begin{ex}
De positieve gehele getallen ($\PosInt$) kun je gewoon opsommen met de
identiteitsfunctie $f(n) = n$. Voor de natuurlijke getallen $\Nat$
werkt de functie $g(n) = n - 1$.
\end{ex}

\begin{ex}
Neem de functies $f\colon \PosInt \to \PosInt$ en $g \colon \PosInt \to
\PosInt$ die zijn gegeven door
\begin{align*}
f(n) & = 2n \text{ en}\\
g(n) & = 2n - 1
\end{align*}
De eerste somt de positieve even gehele getallen op; de tweede doet
hetzelfde voor de positieve oneven gehele getallen. Toch is geen van
beide een opsomming van $\PosInt$, want geen van beide is
!!{surjective}.
\end{ex}

\begin{prob}
Geef een opsomming van de positieve kwadraten $1$, $4$, $9$, $16$, \dots
\end{prob}

\begin{ex}
De functie $f(n) = (-1)^{n} \lceil \frac{(n-1)}{2}\rceil$ gebruikt
$\lceil x \rceil$, de \emph{plafondfunctie}. Die geeft het kleinste
gehele getal dat groter is dan of gelijk is aan $x$. De functie somt
de gehele getallen~$\Int$ op. Je ziet hoe $f$ de getallen uit $\Int$
achtereenvolgens geeft door heen en weer te ``springen'' tussen positieve en negatieve
gehele getallen:
\[
\begin{array}{c c c c c c c c}
f(1) & f(2) & f(3) & f(4) & f(5) & f(6) & f(7) & \dots \\ \\
- \lceil \tfrac{0}{2} \rceil & \lceil \tfrac{1}{2}\rceil & - \lceil \tfrac{2}{2} \rceil & \lceil \tfrac{3}{2} \rceil & - \lceil \tfrac{4}{2} \rceil  & \lceil \tfrac{5}{2}
\rceil & - \lceil \tfrac{6}{2} \rceil & \dots \\ \\
0 & 1 & -1 & 2 & -2 & 3 & \dots
\end{array}
\]
Je kunt $f$ ook zo per geval beschrijven:
\[
f(n) = \begin{cases}
  0 & \text{als $n = 1$}\\
  n/2 & \text{als $n$ even is}\\
  -(n-1)/2 & \text{als $n$ oneven en $>1$ is}
  \end{cases}
\]
\end{ex}

\begin{prob}
  Laat zien dat als $A$ en $B$ !!{enumerable} zijn, ook $A \cup B$ dat
  is. Ga zo te werk: neem aan dat er !!{surjective} functies
  $f\colon \PosInt \to A$ en $g\colon \PosInt \to B$ zijn. Maak een
  !!a{surjective} functie~$h\colon \PosInt \to A \cup B$ en bewijs dat
  die !!{surjective} is. Bekijk ook wat er gebeurt als $A$ leeg is
  of~$B = \emptyset$.
\end{prob}
  
\begin{prob}
  Laat zien dat als $B \subseteq A$ en $A$ !!{enumerable} is, ook~$B$
  dat is. Neem aan dat er een !!a{surjective} functie
  $f\colon \PosInt \to A$ is. Maak een !!a{surjective} functie
  $g\colon \PosInt \to B$ en bewijs dat die !!{surjective} is. Wat
  gebeurt er als $B = \emptyset$?
\end{prob}
    
\begin{prob}
  Laat met inductie naar $n$ zien dat als $A_1$, $A_2$, \dots, $A_n$
  allemaal !!{enumerable} zijn, ook $A_1 \cup \dots \cup A_n$ dat is.
  Je mag aannemen dat als twee verzamelingen $A$ en~$B$ !!{enumerable}
  zijn, ook~$A \cup B$ dat is. 
\end{prob}

Wanneer je de !!{element}s van een verzameling op een lijst zet, ligt
het misschien meer voor de hand om vanaf het $1$-ste !!{element} te
tellen. Maar wiskundigen tellen graag met de natuurlijke getallen~$\Nat$.
Ze noemen de !!{element}s van een lijst daarom het $0$-de, $1$-ste,
$2$-de, enzovoort. We kunnen een opsomming dus net
zo goed definiëren als een !!a{surjective} functie van $\Nat$ naar~$A$.
Beide definities zeggen hetzelfde.

\begin{prop}\ollabel{prop:enum-shift}
  Er bestaat een !!a{surjection} $f\colon \PosInt \to A$ precies
  wanneer er een !!a{surjection} $g\colon \Nat \to A$ bestaat.
\end{prop}

\begin{proof}
  Begin met een !!a{surjection} $f\colon \PosInt \to A$. Definieer
  $g(n) = f(n+1)$ voor elke $n \in \Nat$. Dan is $g\colon \Nat \to A$
  !!{surjective}. Omgekeerd beginnen we met een !!a{surjection}
  $g\colon \Nat \to A$ en definiëren we $f(n) = g(n-1)$.
\end{proof}

Daarom geldt:

\begin{cor}\ollabel{cor:enum-nat}
Een verzameling $A$ is !!{enumerable} als zij leeg is of als er een
!!a{surjective} functie $f\colon \Nat \to A$ bestaat, en alleen dan.
\end{cor}

We zagen al dat je uit een lijst met !!{element}s van een
verzameling~$A$ alle herhalingen kunt halen. Bij opsommingen kan dat
ook, maar we moeten iets zorgvuldiger zeggen en bewijzen wat er dan
gebeurt. Een functie $f\colon \PosInt \to A$ moet voor elke
$n \in \PosInt$ een waarde hebben. Als~$A$ maar eindig veel
!!{element}s heeft, zijn er dus oneindig veel !!{element}s van~$\PosInt$
waarvoor de functie een waarde moet geven, maar slechts eindig veel
mogelijke !!{element}s van~$A$. Dan kun je niet voorkomen dat dezelfde
waarde vaker voorkomt. Is de lijst zonder herhalingen eindig, dan
moeten we voor~$f$ een ander domein nemen dat ook maar eindig veel
!!{element}s heeft. Zonder herhalingen moet $f$ !!{injective} zijn.
De functie is ook !!{surjective}. We zoeken dus een !!a{bijection}
tussen een eindige verzameling $\{1, \dots, n\}$ of $\PosInt$ en~$A$.

\begin{prop}\ollabel{prop:enum-bij}
Als $f\colon \PosInt \to A$ !!{surjective} is en dus een opsomming
van~$A$ is, bestaat er een !!a{bijection} $g\colon Z \to A$. Daarbij
is $Z$ gelijk aan~$\PosInt$ of aan $\{1, \dots, n\}$ voor een
$n \in \PosInt$.
\end{prop}

\begin{proof}
  We definiëren $g$ recursief: elke nieuwe waarde kiezen we met behulp
  van de waarden die al gekozen zijn. Begin met
  $g(1) = f(1)$. Is $g(i)$ al gekozen, dan nemen we voor $g(i+1)$ de
  eerste waarde in $f(1)$, $f(2)$, \dots{} die nog niet voorkomt onder
  $g(1)$, \dots, $g(i)$, als zo'n waarde bestaat. Heeft $A$ precies $n$
  !!{element}s, dan zijn $g(1)$, \dots, $g(n)$ allemaal gekozen. We
  hebben dan een functie $g\colon \{1, \dots, n\} \to A$. Heeft $A$
  oneindig veel !!{element}s, dan moet er voor elke $i$ een
  !!a{element} van~$A$ in de opsomming $f(1)$, $f(2)$, \dots, staan dat
  nog niet voorkomt onder $g(1)$, \dots, $g(i)$. In dat geval hebben we
  een functie $g\colon \PosInt \to A$.
  
  De functie $g$ is !!{surjective}. Elk element van~$A$ komt immers voor
  in de lijst $f(1)$, $f(2)$, \dots{} omdat $f$ !!{surjective} is. Het wordt
  daarom uiteindelijk de waarde van~$g(i)$ voor een bepaalde~$i$.
  De functie is ook !!{injective}. Stel namelijk dat er een $j < i$
  was waarvoor $g(j) = g(i)$. Dan zou $g(i)$ al in de lijst $g(1)$, \dots,
  $g(i-1)$ voorkomen. Dat is in strijd met de manier waarop we~$g$ hebben
  gekozen.
\end{proof}

\begin{cor}\ollabel{cor:enum-nat-bij}
Een verzameling $A$ is !!{enumerable} als zij leeg is of als er een
!!a{bijection} $f\colon N \to A$ bestaat, waarbij $N = \Nat$ of
$N = \{0, \dots, n\}$ voor een $n \in \Nat$, en alleen dan.
\end{cor}

\begin{proof}
$A$ is !!{enumerable} als $A$ leeg is of als er een !!a{surjective}
$f\colon \PosInt \to A$ bestaat, en alleen dan. Uit
\olref{prop:enum-bij} volgt dat dit precies het geval is als er een
!!a{bijective} functie~$f\colon Z \to A$ bestaat waarvoor $Z =
\PosInt$ of $Z = \{1, \dots, n\}$ voor een $n \in \PosInt$. Hetzelfde
argument als in het bewijs van \olref{prop:enum-shift} laat zien dat
dit weer precies het geval is als er een !!a{bijection}
$g\colon N \to A$ bestaat, waarbij $N = \Nat$ of
$N = \{0, \dots, n-1\}$.
\end{proof}

\begin{prob}
  Volgens \olref[sfr][siz][enm]{defn:enumerable} is een verzameling $A$
  aftelbaar als $A = \emptyset$ of als er een !!a{surjective}
  $f\colon \PosInt \to A$ bestaat, en alleen dan. Je kunt een
  ``!!{enumerable} verzameling'' ook op een andere, even precieze manier definiëren: een
  verzameling is aftelbaar als er een !!a{injective} functie
  $g\colon A \to \PosInt$ bestaat, en alleen dan. Laat zien dat beide
  definities hetzelfde zeggen. Laat dus zien dat er een !!a{injective}
  functie $g\colon A \to \PosInt$ bestaat precies wanneer
  $A = \emptyset$ of er een !!a{surjective} $f\colon \PosInt \to A$
  bestaat.
\end{prob}
\end{document}
